% !PRINTED 589
\weberpage{589}
%<<CONT

The product of these roots is therefore, since $\boldsymbol{s}^m$ has the coefficient 1,
equal to the negative independent term, i.\,e. it is:
\weq{22}{\varepsilon\,\frac{\boldsymbol{s}(\mu u)}{\boldsymbol{s}(u)} = \overset{\varrho}{\prod} \boldsymbol{s}\left(u + \frac{2\varrho}{\mu}\right),}
and from this, for $u = 0$, as above:
\weq{23}{\varepsilon\mu = \overset{\varrho}{\prod} \boldsymbol{s}\left(\frac{2\varrho}{\mu}\right).}

Now let $\eta$ be a fractional number in $\Omega$ with the odd
ideal denominator $\mathfrak{a}$; then it follows from (22) that $\boldsymbol{s}(\mu\eta)/\boldsymbol{s}(\eta)$ is an
integer, if $\mu$ is an arbitrary number (6) in $\Omega$. If we take
$\mu$ relatively prime to $\mathfrak{a}$ and set $\mu'\mu \equiv 1$ (mod $\mathfrak{a}$), then we may
replace $\eta$ by $\mu'\eta$ and find that $\boldsymbol{s}(\eta)/\boldsymbol{s}(\mu'\eta)$
is likewise an integer, in which $\mu'$ is just as arbitrary as $\mu$. Consequently
$\boldsymbol{s}(\eta)/\boldsymbol{s}(\mu\eta)$ is also an integer and hence a \gesperrt{unit}.

We have therefore the theorem:

3. \gesperrt{If $\eta$ is a number in $\Omega$ with odd ideal denominator,
and if $\mu$ runs through a sequence of numbers relatively prime
to $\mathfrak{a}$, then the numbers}
\[
\boldsymbol{s}(\mu\eta)
\]
\gesperrt{are associated with one another.}

Yet a further consequence follows from this if one sets
$u = 2\varrho_1/\mu_1$ in (22):
\weq{24}{\frac{\boldsymbol{s}\left(\dfrac{2\mu\varrho_1}{\mu_1}\right)}{\boldsymbol{s}\left(\dfrac{2\varrho_1}{\mu_1}\right)} = \overset{\varrho}{\prod} \boldsymbol{s}\left(\frac{2\varrho_1}{\mu_1} + \frac{2\varrho}{\mu}\right).}

If now $\mu$ and $\varrho_1$ are relatively prime to $\mu_1$, then by 3. the left side
is a unit, and consequently all the factors of the
right side
\weq{25}{\boldsymbol{s}\left(\frac{2\varrho_1}{\mu_1} + \frac{2\varrho}{\mu_2}\right),}
% [source sic: equation (25) prints \mu_2, whereas equation (24) and the context use \mu]
which are indeed integers, must be units.

But every $\boldsymbol{s}(\eta)$ can be brought into this form if
\[
\eta = \frac{\nu}{\mu}
\]
is a fractional number in $\Omega$ whose denominator is odd and divisible by
two distinct prime ideals. For in this case one
can decompose the denominator $\mu$ of $\eta$ into two ideals $\mathfrak{a}_1\mathfrak{a}_2$
%CONT>>
